\documentclass[10pt,a4paper]{amsart}
\usepackage{amsmath,amssymb,amsthm,graphicx}
% Retained minimal preamble for the Stacks Project chapter flat.tex (More on Flatness).
% Verbatim excerpt of the chapter preamble: Xy-pic setup (native preamble lines 14-18) and the native
% theorem declarations with the mathematical macros (native preamble lines 154-217). Copied mechanically;
% nothing added, removed or redefined. The wrapper class and the amsmath/amssymb/amsthm/graphicx base
% packages are supplied by the assembler, so the class/externaldocument part of the native preamble is omitted.
% For commutative diagrams we use Xy-pic
\usepackage[all]{xy}

% We use 2cell for 2-commutative diagrams.
\xyoption{2cell}
% Theorem environments.
%
\theoremstyle{plain}
\newtheorem{theorem}[subsection]{Theorem}
\newtheorem{proposition}[subsection]{Proposition}
\newtheorem{lemma}[subsection]{Lemma}

\theoremstyle{definition}
\newtheorem{definition}[subsection]{Definition}
\newtheorem{example}[subsection]{Example}
\newtheorem{exercise}[subsection]{Exercise}
\newtheorem{situation}[subsection]{Situation}

\theoremstyle{remark}
\newtheorem{remark}[subsection]{Remark}
\newtheorem{remarks}[subsection]{Remarks}

\numberwithin{equation}{subsection}

% Macros
%
\def\lim{\mathop{\mathrm{lim}}\nolimits}
\def\colim{\mathop{\mathrm{colim}}\nolimits}
\def\Spec{\mathop{\mathrm{Spec}}}
\def\Hom{\mathop{\mathrm{Hom}}\nolimits}
\def\Ext{\mathop{\mathrm{Ext}}\nolimits}
\def\SheafHom{\mathop{\mathcal{H}\!\mathit{om}}\nolimits}
\def\SheafExt{\mathop{\mathcal{E}\!\mathit{xt}}\nolimits}
\def\Sch{\mathit{Sch}}
\def\Mor{\mathop{\mathrm{Mor}}\nolimits}
\def\Ob{\mathop{\mathrm{Ob}}\nolimits}
\def\Sh{\mathop{\mathit{Sh}}\nolimits}
\def\NL{\mathop{N\!L}\nolimits}
\def\CH{\mathop{\mathrm{CH}}\nolimits}
\def\proetale{{pro\text{-}\acute{e}tale}}
\def\etale{{\acute{e}tale}}
\def\QCoh{\mathit{QCoh}}
\def\Ker{\mathop{\mathrm{Ker}}}
\def\Im{\mathop{\mathrm{Im}}}
\def\Coker{\mathop{\mathrm{Coker}}}
\def\Coim{\mathop{\mathrm{Coim}}}

% Boxtimes
%
\DeclareMathSymbol{\boxtimes}{\mathbin}{AMSa}{"02}

%
% Macros for moduli stacks/spaces
%
\def\QCohstack{\mathcal{QC}\!\mathit{oh}}
\def\Cohstack{\mathcal{C}\!\mathit{oh}}
\def\Spacesstack{\mathcal{S}\!\mathit{paces}}
\def\Quotfunctor{\mathrm{Quot}}
\def\Hilbfunctor{\mathrm{Hilb}}
\def\Curvesstack{\mathcal{C}\!\mathit{urves}}
\def\Polarizedstack{\mathcal{P}\!\mathit{olarized}}
\def\Complexesstack{\mathcal{C}\!\mathit{omplexes}}
% \Pic is the operator that assigns to X its picard group, usage \Pic(X)
% \Picardstack_{X/B} denotes the Picard stack of X over B
% \Picardfunctor_{X/B} denotes the Picard functor of X over B
\def\Pic{\mathop{\mathrm{Pic}}\nolimits}
\def\Picardstack{\mathcal{P}\!\mathit{ic}}
\def\Picardfunctor{\mathrm{Pic}}
\def\Deformationcategory{\mathcal{D}\!\mathit{ef}}

% --- appended 2026-10-12 (Chief): the project own definitions of the reference/slogan/history
% display environments, copied mechanically from the shared Stacks preamble lines 8-12
% (\usepackage{verbatim} plus the three \newenvironment lines, which the chapter preamble keeps
% verbatim from the project preamble). Nothing added, removed or redefined: these environments
% carry no mathematical content and the project itself typesets them as comments. Without these
% three lines a counted statement whose own text uses them fails with Environment slogan undefined.
\usepackage{verbatim}
\newenvironment{reference}{\comment}{\endcomment}
\newenvironment{slogan}{\comment}{\endcomment}
\newenvironment{history}{\comment}{\endcomment}
\usepackage{cleveref}
\crefname{originallemma}{original lemma}{original lemmas}
\Crefname{originallemma}{Original lemma}{Original lemmas}
\makeatletter
\newlabel{lemma-homothety-spectrum}{{05DE}{source}{}{sourcepointer-1}{}}
\newlabel{lemma-homothety-spectrum@cref}{{[originallemma][0][]05DE}{[1][0][]source}{}{}{}}
\makeatother
\crefname{originallemma}{original lemma}{original lemmas}
\Crefname{originallemma}{Original lemma}{Original lemmas}
\makeatletter
\newlabel{lemma-invert-universally-injective}{{05FR}{source}{}{sourcepointer-2}{}}
\newlabel{lemma-invert-universally-injective@cref}{{[originallemma][0][]05FR}{[1][0][]source}{}{}{}}
\makeatother
\title{Associated primes of a fibre of a projective module do not meet the extended maximal ideal}
\author{The Stacks Project Authors}
\date{stacks-project:flat.tex:local-snapshot-20261011; GNU FDL 1.2 or later}
\begin{document}
\maketitle
\noindent\textit{Editorial note.} All delivered mathematical text is the author's own, copied from the pinned Stacks Project chapter file flat.tex; no statement or proof step was rewritten, reordered or removed.\par\medskip
\noindent\textit{Editorial note.} Preamble provenance: the chapter's own retained minimal preamble (sources/stacks-flat/retained-preamble.tex, lines 1-85) is a verbatim excerpt of the native Xy-pic setup and of the native theorem declarations and mathematical macros, followed by the project's own reference/slogan/history display-environment definitions copied mechanically from the shared preamble; the wrapper class and the amsmath/amssymb/amsthm/graphicx base packages are supplied by the assembler, and extra\_packages is left empty because everything needed is already in the retained preamble. No macro was invented or redefined.\par\medskip
\noindent\textit{Editorial note.} Same-chapter pointers are located at their \textbackslash{}label lines inside the pinned chapter file; cross-chapter pointers are located in the project's own tag list staged in the snapshot (sources/stacks-chow/tags.txt), at the verbatim lines carrying their canonical tags.\par\medskip
\noindent\textit{Allowed standard prerequisites (editorial).}\begin{itemize}
\item The homothety spectrum and universal injectivity of the localisation away from the zerodivisors of the fibre. Same chapter flat.tex, labels lemma-homothety-spectrum and lemma-invert-universally-injective (proof lines 4331-4332); Algebra, definition algebra-definition-ord for the order used later. Gives the semi-locality statement and the universal injectivity of the map localising away from the elements that are not zerodivisors on the fibre, whose conclusions the proof uses.
\end{itemize}\medskip
\ifdefined\hypertarget\hypertarget{sourcepointer-1}{}\fi
\noindent\textit{Source reference.} Original lemma 05DE: sources/stacks-flat/pinned-original-source.tex, source lines 1783--1783. Navigation pointer for \textbackslash{}ref\{lemma-homothety-spectrum\} used inside the counted span (line 4331); the argument that consumes it is retained in the span.\par\medskip
\ifdefined\hypertarget\hypertarget{sourcepointer-2}{}\fi
\noindent\textit{Source reference.} Original lemma 05FR: sources/stacks-flat/pinned-original-source.tex, source lines 1976--1976. Navigation pointer for \textbackslash{}ref\{lemma-invert-universally-injective\} used inside the counted span (line 4332); the argument that consumes it is retained in the span.\par\medskip
\begin{lemma}
\label{lemma-explain-why-pure}
Let $R$ be a local ring with maximal ideal $\mathfrak m$.
Let $R \to S$ be a ring map. Let $N$ be an $S$-module.
Assume
\begin{enumerate}
\item $N$ is projective as an $R$-module, and
\item $S/\mathfrak mS$ is Noetherian and $N/\mathfrak mN$ is a finite
$S/\mathfrak mS$-module.
\end{enumerate}
Then for any prime $\mathfrak q \subset S$ which is an associated prime of
$N \otimes_R \kappa(\mathfrak p)$ where $\mathfrak p = R \cap \mathfrak q$
we have $\mathfrak q + \mathfrak m S \not = S$.
\end{lemma}

\begin{proof}
Note that the hypotheses of
Lemmas \ref{lemma-homothety-spectrum} and
\ref{lemma-invert-universally-injective}
are satisfied. We will use the conclusions of these lemmas without further
mention. Let $\Sigma \subset S$ be the multiplicative set of elements
which are not zerodivisors on $N/\mathfrak mN$. The map
$N \to \Sigma^{-1}N$ is $R$-universally injective. Hence we see that
any $\mathfrak q \subset S$ which is an associated prime of
$N \otimes_R \kappa(\mathfrak p)$ is also an associated prime of
$\Sigma^{-1}N \otimes_R \kappa(\mathfrak p)$. Clearly this implies that
$\mathfrak q$ corresponds to a prime of $\Sigma^{-1}S$.
Thus $\mathfrak q \subset \mathfrak q'$ where $\mathfrak q'$
corresponds to an associated prime of $N/\mathfrak mN$ and we win.
\end{proof}

\end{document}
